% results.tex (국문 v15, 2026-10-21): 영문 v15 의 쌍둥이. 서론의 질문 순서를 따르는 다섯 절, 수치 뒤에 정성 해석, 외부 CALM 지역 제외. Δ = RMSE(앞) − RMSE(뒤), 같은 지점-연도; 대괄호는 95 % 짝지은 블록 부트스트랩 구간.
\section{결과}\label{sec:results}

\subsection{기후 입력, 검증 설계, 지역 안과 밖의 대비}
먼저 관측 연도의 기후가 중요한지, 각 방법이 보지 않은 지점에서 얼마나 정확한지, 그리고 그 답이 자료를 나누는 방식에 얼마나 좌우되는지를 묻는다. 2011--2024년 평균 기후 입력을 관측 연도의 입력으로 바꾸면 공간 블록 RMSE 는 재보정 Stefan 모델에서 레나 델타($\Delta$RMSE \ndforcLePo{} cm), 물리 잔차 모델에서 레나 델타(\ndforcLeRo{} cm)와 캐나다(\ndforcCaRo{} cm)에서 줄었고, 라벨의 절반이 한 해에서 나온 알래스카에서는 변화가 판정되지 않았다(그림~\ref{fig:accuracy}a). 정확도를 떠나 연도별 입력은 각 예측을 명시된 해에 묶어 준다. 이 입력을 쓴 공간 블록 교차검증에서 물리 잔차 모델은 재보정 Stefan 모델의 RMSE 를 레나 델타에서 \nrmseLeblkPo{} cm 에서 \nrmseLeblkRo{} cm 로($\Delta$RMSE \ndrefLeblkRo{} cm), 캐나다에서 \nrmseCablkPo{} cm 에서 \nrmseCablkRo{} cm 로(\ndrefCablkRo{} cm) 줄였고, 알래스카에서는 두 모델이 \nrmseAkblkPo{} cm 와 \nrmseAkblkRo{} cm 에 이르렀다(\ndrefAkblkRo{} cm; 그림~\ref{fig:accuracy}b). 관측 대 예측 산점도는 학습기가 무엇을 더하는지 보여 준다. 재보정 Stefan 모델은 어느 지역에서나 좁은 범위의 값만 예측하는데, 각 해의 기후에 계수 하나를 적용해서는 지점들을 가를 수 없기 때문이다. 물리 잔차 모델은 레나 델타와 캐나다에서 예측을 1:1 선을 따라 퍼뜨린다(보충 그림~S11). 응답 전체를 라벨에서 배우는 직접 ML 은 한 지역의 라벨을 모두 쓸 때 물리 잔차 모델보다 낮은 점 오차에 이르렀으나(알래스카 \nrmseAkblkDz{} 대 \nrmseAkblkRo, 레나 델타 \nrmseLeblkDz{} 대 \nrmseLeblkRo, 캐나다 \nrmseCablkDz{} 대 \nrmseCablkRo{} cm), 재보정 Stefan 모델과의 차이는 캐나다에서만 판정되었다(\ndrefCablkDz{} cm). 같은 예측을 지점-연도 대신 1 km 셀-연도 평균으로 채점하면 직접 ML 의 우위는 알래스카에서 사라지고(\ncyAkblkDz{} 대 \ncyAkblkRo{} cm) 다른 지역에서도 좁아진다(보충 표~S6). 이는 그 우위의 일부가 셀 안 탐침 지점 사이의 차이를 맞추는 데서 오며, 1 km 지도는 그 정보를 쓸 수 없음을 뜻한다.

\paperfig{Fig3_accuracy}{1.0}{\textbf{기후 입력, 검증 설계, 지역 안과 밖의 대비.} \textbf{a}, 2011--2024년 평균 기후 입력을 관측 연도의 입력으로 바꿨을 때의 공간 블록 RMSE 변화(95\% 짝지은 블록 부트스트랩 구간; 음수는 연도별 입력이 유리함). \textbf{b}, 덜 엄격한 설계부터 가장 엄격한 설계까지 다섯 검증 설계에서 세 방법의 RMSE. 칸의 색은 같은 설계의 재보정 Stefan 모델과의 차이(파랑, 낮음; 빨강, 높음)이고, 지역 홀드아웃에서는 다른 두 지역의 계수를 쓴 Stefan 모델이 기준이다. \textbf{c}, 그 지역의 라벨을 쓸 때(공간 블록 교차검증)와 그 지역을 빼고 다른 두 지역으로만 학습했을 때의 같은 방법들의 RMSE.}{fig:accuracy}

방법의 순위는 검증 설계에 따라 달랐다(그림~\ref{fig:accuracy}b, 보충 표~S4). 시험 지점을 학습 지점 옆에 두는 무작위 교차검증은 직접 ML 에 가장 낮은 오차를 주었고(알래스카 \nrmseAkrandDz{} cm), 공간 블록 교차검증은 직접 ML 의 오차를 알래스카 · 레나 델타 · 캐나다에서 \noptimAkDz, \noptimLeDz, \noptimCaDz{} cm 높였으며, 연도 홀드아웃과 블록-연도 설계는 이를 더 높였다. 한 지역 전체를 빼면 레나 델타의 직접 ML 오차는 \nrmseLeregDz{} cm 로 올랐고 다른 지역의 계수를 쓴 Stefan 모델은 \nrmseLeregPz{} cm 였다($\Delta$RMSE \ndrefLeregDz{} cm). 물리 잔차 모델은 그곳에서 Stefan 모델과 \ndrefLeregRoabs{} cm 안에 머물렀고 캐나다에서는 이를 개선하였다(\ndrefCaregRo{} cm; 그림~\ref{fig:accuracy}c). 즉 시험 지점이 학습 지점에서 멀어질수록 세 방법은 다르게 행동한다. Stefan 모델은 계수 하나가 과적합할 수 없어 설계에 거의 둔감하고, 직접 ML 은 시험 지점과 학습 지점이 섞여 있을 때 가장 정확하고 그렇지 않을 때 가장 부정확하며, 물리 잔차 모델은 지역 안 이득의 대부분을 지키면서 기준값의 안정성을 물려받는다. 지도는 관측이 없는 셀과 해에 필요하므로, 이하에서는 공간 블록 교차검증을 주 설계로, 지역 홀드아웃을 전이성 시험으로 쓴다.

\subsection{물리 잔차 이득의 조건}
다음으로 학습기에 어떤 정보가 필요하고 보정이 어디서 작용하는지를 묻는다. 이득은 잔차 학습기가 쓸 수 있는 공변량에 좌우되었다(그림~\ref{fig:conditions}a). 기후 공변량만으로는 잔차 모델이 알래스카(\ndfeatAkFclim{} cm)와 캐나다(\ndfeatCaFclim{} cm)에서 Stefan 모델을 개선하지 못하였고, 지형과 토양 공변량을 더하자 위의 이득이 나타났다. MODIS 식생 · 적설 시기나 고해상 토양 · 토지피복 분류는 레나 델타와 캐나다에서 결과를 0.5 cm 미만으로 바꿨고, 후자는 알래스카에서 작은 이득을 더하였다(\ndfeatAkFhires{} cm). 공변량 여덟 개의 지역 내 3분위에 걸쳐 이득은 레나 델타의 대부분 조건(구간이 0 을 제외한 3분위 \ngcLesig{} / \ngcLecells{})과 캐나다(\ngcCasig{} / \ngcCacells)에서 나타났고 Stefan 모델의 오차와 함께 커졌으며(3분위 간 상관 레나 델타 \ngcLecorr, 캐나다 \ngcCacorr; 그림~\ref{fig:conditions}b), 레나 델타의 가장 따뜻한 TDD 3분위에서 \ndtddLewarm{} cm 에 이르렀다. 알래스카에서는 어느 3분위도 \ngcAkmaxabs{} cm 넘게 바뀌지 않았다. 이득의 부재는 특정 조건에 국한되지 않고 고르며, 알래스카 자료의 두 특징과 부합한다. 지역 계수가 이미 지점들을 잘 나타내고(Stefan RMSE \nrmseAkblkPo{} cm, 관측 범위 대비 세 지역 중 최저), 남은 오차는 1 km 셀 안 탐침 지점 사이의 산포(합동 셀 내 표준편차 \ncyWithinSdAk{} cm)가 지배하는데 이는 어떤 1 km 공변량도 해소하지 못한다. 잔차 모델의 평균 절대 SHAP 값에서 지형 변수는 레나 델타에서 \nshapLeResTerrain\%, 토양 변수는 캐나다에서 \nshapCaResSoil\% 를 차지하였다(보충 그림~S3). 물리를 다른 곳에 두거나(유사라벨은 Stefan 모델을 재현, 레나 델타 \ndrefLeblkDo{} cm; 입력 특징은 직접 ML 처럼 행동) CatBoost 를 다른 학습기 넷으로 바꿔도 레나 델타와 캐나다의 이득 부호는 바뀌지 않았다(보충 그림~S2).

\paperfig{Fig4_conditions}{1.0}{\textbf{물리 잔차 이득의 조건과 라벨이 적은 지역으로의 전이.} \textbf{a}, 여섯 겹 공변량 묶음에 대한 물리 잔차 모델과 재보정 Stefan 모델의 공간 블록 RMSE 차이. \textbf{b}, 공변량 여덟 개의 지역 내 3분위에서의 같은 차이. 테두리 칸은 구간이 0 을 제외하고, 빗금 칸은 지점-연도 30개 또는 블록 5개 미만이다. \textbf{c}, 다른 두 지역과 대상 지역 라벨 $n$개로 학습한 모델의 지역별 RMSE. 축을 벗어난 값은 숫자로 적었다. a 와 b 의 구간은 95\% 짝지은 블록 부트스트랩 구간이고, 음수는 물리 잔차 모델이 유리함을 뜻한다.}{fig:conditions}

\subsection{지역 안의 라벨 수와 새 지역으로의 전이}
세 번째 질문은 학습기에 지역 관측이 몇 개나 필요한가이다. 한 지역의 지점-연도에서 뽑은 학습 라벨 25개로 물리 잔차 모델은 같은 라벨로 재보정한 Stefan 모델만큼은 정확하였고(알래스카 \nllAkRotwfive, 레나 델타 \nllLeRotwfive, 캐나다 \nllCaRotwfive{} cm), 라벨 100개에서는 세 지역 모두에서 이득이 판정되었다(\nllAkRohundred, \nllLeRohundred, \nllCaRohundred{} cm; 보충 그림~S4). 라벨 25개로 학습한 직접 ML 은 알래스카(\nllAkDztwfive{} cm)와 캐나다(\nllCaDztwfive{} cm)에서 Stefan 모델보다 점 오차가 높았고, 알래스카와 레나 델타에서 수백 개의 라벨에서만 물리 잔차 모델을 앞섰다. 두 곡선은 서로 모순되지 않는다. 라벨이 수천 개면 직접 ML 이 더 가깝게 맞추지만, 물리 잔차 모델은 그 일부의 라벨로 거의 같은 정확도에 이르고 기준값 아래로 떨어지지 않는다. 다른 모델 구조는 보충 그림~S4 에 있다.

전이 실험에서는 각 지역을 다른 두 지역과 대상 지역 라벨 $n$개로 학습한 모델로 예측하였다(그림~\ref{fig:conditions}c; $n \leq 500$, 대상 지역의 학습 겹에서 추출). 대상 라벨 없이 직접 ML 은 레나 델타에서 \ntlrmseLeDszero{} cm 였고, 원천 계수 Stefan 모델은 \ntlrmseLePzzero{} cm, 물리 잔차 모델은 \ntlrmseLeRszero{} cm 였다. 대상 라벨 25개는 물리 잔차 오차를 \ntlrmseLeRstwfive{} cm 로, 직접 ML 오차를 \ntlrmseLeDstwfive{} cm 로 줄였다. 알래스카에서는 물리 잔차 모델이 모든 $n$ 에서 가장 낮은 오차를 보였다(대상 라벨 없이 \ntlrmseAkRszero{} cm). 캐나다에서는 원천 지역으로 학습한 직접 ML 이 가장 정확하였다(\ntlrmseCaDszero{} cm). 캐나다 라벨 25개로 Stefan 계수를 재보정하면 오차가 \ntlrmseCaPzzero{} cm 에서 \ntlrmseCaPntwfive{} cm 로 늘었고, 이 기준값을 쓰는 물리 잔차 모델도 뒤따랐다(\ntlrmseCaRszero{} 에서 \ntlrmseCaRstwfive{} cm). 이는 캐나다 지점의 융해 깊이 범위가 넓어(탐침 측선 관측 평균 \nsrcCaABoVEObs{} cm, CALM 격자 \nsrcCaCALMObs{} cm) 라벨 25개로 맞춘 계수 하나가 이를 잘 나타내지 못하기 때문인 것으로 추정된다.

\subsection{같은 셀 · 같은 해에서의 공개 제품 비교}
다음으로 공개 제품이 우리 라벨이 있는 곳에서 우리 모델과 같은 방식으로 채점했을 때 어떤 성능을 보이는지 묻는다. 레나 델타의 공통 지점-연도 \nccLeN{}개(2010--2020년; Wei 등, Liu 등, ESA CCI v4 모두 관측 연도 값을 제공하는 지점-연도)에서 물리 잔차 모델의 RMSE 는 \nccLeRoall{} cm 였고, ESA CCI v5 는 \nccLeCcifiveall{} cm, Aalto 등은 \nccLeAaltoall{} cm, ESA CCI v4 는 \nccLeCcifourall{} cm, Liu 등은 \nccLeLiuall{} cm, Wei 등은 \nccLeWeiall{} cm 였다(그림~\ref{fig:products}a, 보충 표~S5). Wei 등의 학습 지점에서 25 km 넘게 떨어진 레나 델타 지점-연도(Wei 등 값이 있는 \nnprodLeweimask{}개, ESA CCI v5 값이 있는 \nnprodLeccifivemask{}개)에서 짝지은 차이는 Wei 등 대비 \ndprodLeweiRomask{} cm, ESA CCI v5 대비 \ndprodLeccifiveRomask{} cm 였다. 알래스카(\nccAkRoall{} 대 Wei 등 \nccAkWeiall, Liu 등 \nccAkLiuall{} cm)와 캐나다(\nccCaRoall{} 대 \nccCaWeiall, \nccCaLiuall{} cm)에서는 모델과 두 ML 제품이 동급이었다.

관측 대 예측 산점도는 지역이 왜 다른지 보여 준다(그림~\ref{fig:products}b). 레나 델타에서 두 ML 제품은 관측 ALT 와 무관하게 대부분의 지점에 약 40--70 cm 의 값을 준다. 관측과 예측의 상관은 Wei 등 \nrLeWei, Liu 등 \nrLeLiu{} 로 물리 잔차 모델의 \nrLeRo{} 와 대비되고, 평균 편향은 \nbiasLeAllWei{} cm 와 \nbiasLeAllLiu{} cm 로 물리 잔차 모델의 \nbiasLeAllRo{} cm 와 대비되며, 과대 추정은 델타 북부의 얇은 활동층에서 가장 컸다(북위 72$^\circ$ 이북 \nbiasLeNorthWei{} cm, \nbiasLeNorthLiu{} cm). 오차 분포(그림~\ref{fig:products}c)는 이 체계적 편향을 산포와 분리한다. 레나 델타에서 제품은 약 20 cm 치우쳐 있으나 산포는 우리 모델과 비슷하고, 알래스카와 캐나다에서는 모든 계열이 0 근처에 중심을 두고 주로 산포에서 다르다. 이 지역 간 대비는 제품의 학습 자료를 따른다. 알래스카 라벨 지점은 Wei 등의 학습 지점에서 중앙값 \ndistweiAkmed{} km 떨어져 있고, Liu 등은 우리가 쓴 알래스카 · 캐나다 관측의 이전 판인 ABoVE 활동층 자료로 모델을 학습하였으므로 그 지점에서의 제품 값은 우리 라벨과 독립이 아니다. 레나 델타 지점은 Wei 등의 학습 지점에서 중앙값 \ndistweiLemed{} km 떨어져 있다(보충 그림~S6). 레나 델타 라벨 없이 학습한 모델이 이미 같은 지점-연도에서 Wei 등보다 정확하였으므로(\ntlvprmseLeRsWeizero{} 대 \ntlvprmseLeWeizero{} cm; $\Delta$RMSE \ntlvpLeRsWeizero{} cm), 이 지역의 우위는 지역 라벨에만 의존하지 않는다. 연도별로 물리 잔차 모델은 지점-연도 30개 이상인 레나 델타의 \nywLen{}개 해 모두와 알래스카 \nywAkRoWein{}개 해 중 \nywAkRoWeiwin{}개에서 Wei 등보다 오차가 낮았다(보충 그림~S5).

\paperfig{Fig5_products}{1.0}{\textbf{같은 지점 · 같은 해에서의 공개 제품 비교.} \textbf{a}, 공통 지점-연도(2010--2020년)에서 재보정 Stefan 모델, 물리 잔차 모델, 제품 5종의 RMSE. 채운 표지는 전체 지점, 빈 표지는 Wei 등의 학습 지점에서 25 km 넘게 떨어진 지점이며, 45 cm 를 넘는 값은 압축한 구간에 그렸다. \textbf{b}, 공통 지점-연도에서 물리 잔차 모델, Wei 등, Liu 등의 관측 대 예측 ALT 와 RMSE · 평균 편향 · 상관. \textbf{c}, 같은 지점-연도에서 예측 − 관측의 분포(바이올린은 $\pm$80 cm 에서 자름; 점, 중앙값; 막대, 사분위 범위). 우리 모델은 모두 공간 블록 겹 밖 예측이다.}{fig:products}

지역별로 관측 셀이 가장 많은 해(알래스카 · 캐나다 2013년, 레나 델타 2016년)의 지도는 이 차이를 공간적으로 보여 준다(그림~\ref{fig:prodmaps}). 알래스카 관측 셀 \nfeAkN{}개에서 세 지도의 RMSE 는 비슷하였다(\nfeAkOof, \nfeAkWei, \nfeAkLiu{} cm). 관측과의 상관은 우리 모델(\nfeAkROof)이 제품(\nfeAkRWei, \nfeAkRLiu)보다 낮았으나, 이 중 \nfeAkNNear{}개 셀은 Wei 등의 학습 지점에서 25 km 안에 있고, 그 밖의 \nfeAkNFar{}개 셀에서는 세 지도의 오차가 비슷하였다(\nfeAkFarOof, \nfeAkFarWei, \nfeAkFarLiu{} cm; $r$ \nfeAkFarROof, \nfeAkFarRWei, \nfeAkFarRLiu). 지도 간 차이는 고도 600 m 이상의 산지에서 가장 컸고, 북위 68--70$^\circ$의 산지에서 제품은 우리 지도보다 \nmapAkFootMtnWei{} cm(Wei 등), \nmapAkFootMtnLiu{} cm(Liu 등) 깊었다. 레나 델타(2016년 관측 셀 \nfeLeN{}개)에서는 두 제품이 관측보다 깊었다(평균 편향 \nfeLeBiasWei{} cm, \nfeLeBiasLiu{} cm; RMSE \nfeLeWei{} cm, \nfeLeLiu{} cm 대 \nfeLeOof{} cm). 캐나다(북서부)에서는 관측 셀 \nfeCaN{}개의 RMSE 가 우리 모델 \nfeCaOof{} cm, Wei 등 \nfeCaWei{} cm(평균 편향 \nfeCaBiasWei{} cm), Liu 등 \nfeCaLiu{} cm였고, Liu 등의 학습 자료에는 이 관측의 이전 판이 포함된다. Wei 등의 학습 지점에서 25 km 밖의 \nfeCaNFar{}개 셀에서는 \nfeCaFarOof, \nfeCaFarWei, \nfeCaFarLiu{} cm였다.

\paperfig{Fig6_product_maps}{1.0}{\textbf{지역별로 관측 셀이 가장 많은 해의 물리 잔차 모델 지도와 공개 제품 지도, 그리고 관측값.} \textbf{a--c}, 알래스카 2013년의 물리 잔차 모델, Wei 등, Liu 등 지도. 제품 지도는 우리 지도의 셀에만 표시하고, 관측 셀은 관측 ALT를 같은 색 눈금으로 채운 원으로 그렸다. \textbf{d}, 세 지도 모두 값이 있는 관측 셀 \nfeAkN{}개의 관측 대 예측 ALT(1:1 선). 채운 표지는 Wei 등의 학습 지점에서 25 km 안의 셀, 빈 표지는 그 밖의 셀이며, RMSE 는 패널 아래에 전체 셀과 빈 표지 셀에 대해 적었다. 우리 예측은 공간 블록 겹 밖 값이다. \textbf{e--h}, 레나 델타 2016년의 같은 비교(\nfeLeN{}개 셀). \textbf{i--l}, 캐나다(북서부) 2013년의 같은 비교(\nfeCaN{}개 셀). 지역마다 적용 가능 영역 안 세 지도 값의 5--95 백분위수를 색 범위로 쓰며, 옅은 색은 영역 밖의 셀이다.}{fig:prodmaps}

잔차 학습기는 제품을 기준값으로 쓸 수도 있다(식~(\ref{eq:prod})). 레나 델타 라벨로는 Stefan 기준 모델이 제품 기준 모델보다 정확하였고(\npadLeRoRp{} cm; 보충 그림~S7), 알래스카와 캐나다에서는 두 기준값이 구별되지 않았다.

\subsection{2011--2024년 연도별 1 km ALT 지도}
마지막으로 각 지역의 모든 지점-연도로 적합한 물리 잔차 모델을 2011년부터 2024년까지 각 해의 ERA5-Land 기후로 적용하였다(그림~\ref{fig:maps}). 지도는 알래스카 \nncellmapAk{}개, 레나 델타 \nncellmapLe{}개, 캐나다(북서부) \nncellmapCa{}개 셀을 덮는다. 2011--2024년 평균 대비 지역 평균 편차는 알래스카에서 \nanminyrAk{}년 \nanminAk{} cm 에서 \nanmaxyrAk{}년 \nanmaxAk{} cm, 레나 델타에서 \nanminyrLe{}년 \nanminLe{} cm 에서 \nanmaxyrLe{}년 \nanmaxLe{} cm, 캐나다에서 \nanminyrCa{}년 \nanminCa{} cm 에서 \nanmaxyrCa{}년 \nanmaxCa{} cm 였다(그림~\ref{fig:maps}f). 이 편차는 각 해의 기후에 대한 모델의 응답이다. 지역 평균 편차는 9월 1일까지의 지역 평균 TDD 를 알래스카 · 레나 델타 · 캐나다에서 상관 \ntddanomAk, \ntddanomLe, \ntddanomCa{} 로 따르므로(그림~\ref{fig:maps}d), 지도의 가장 깊은 지역 융해는 기간 중 가장 따뜻한 융해 계절(\ntddmaxAk, \ntddmaxLe, \ntddmaxCa{}년)에, 가장 얕은 융해는 가장 추운 계절(\ntddminAk, \ntddminLe, \ntddminCa{}년)에 놓인다. 잔차 학습기는 이 응답을 공간적으로 조절하며, 그래서 더운 해의 편차는 알래스카 내륙 저지에서 가장 크고 노스슬로프에서 가장 작다(그림~\ref{fig:maps}f). 관측 셀의 연도별 평균과의 비교는 캠페인마다 측정 지점 집합이 바뀌어 제한되므로, 보충 그림~S10 에서 장기 관측 지점의 지도 시계열과 관측 시계열을 대신 비교한다. 연도별 지도의 선형 추세는 알래스카 · 레나 델타 · 캐나다 셀에서 평균 \ntrendAk, \ntrendLe, \ntrendCa{} cm/10년이었고 양수 셀은 \ntrendposAk\%, \ntrendposLe\%, \ntrendposCa\% 였다(그림~\ref{fig:maps}e). 이 추세는 모델이 전달한 그 기간의 ERA5-Land 기후를 나타내며 관측 추세가 아니다. 적용 가능 영역은 알래스카 셀의 \naoaAk\%, 레나 델타 셀의 \naoaLe\%, 캐나다 셀의 \naoaCa\% 를 덮고, 영역 밖 셀은 주로 산지와 영구동토 남쪽 경계에 있다(보충 그림~S8). 90\% 예측 구간의 반폭은 \nhwAk, \nhwLe, \nhwCa{} cm 이다. 잔차 보정은 보충 그림~S9 에 있다.

\paperfig{Fig7_maps}{1.0}{\textbf{2011--2024년 연도별 1 km ALT 지도.} \textbf{a--c}, 알래스카, 레나 델타, 캐나다(북서부) 연도별 지도의 2011--2024년 평균 ALT 와 오른쪽의 확대 창(1 km 셀), 관측 평균 ALT 로 채운 관측 셀. 옅은 색은 적용 가능 영역 밖의 셀이다. \textbf{d}, 연도별 지도의 지역 평균 ALT 편차와 같은 해의 지역 평균 TDD 편차(해마다 점 하나, 최소제곱 직선; 가장 따뜻한 해와 가장 추운 해를 표시). \textbf{e}, 연도별 지도의 선형 추세(셀별 최소제곱). 각 해의 기후에 대한 모델 응답을 나타낸다. \textbf{f}, 알래스카의 2011--2024년 평균 대비 연도별 ALT 편차.}{fig:maps}
